% from ega0-1-fr.tex
\oldpage[0\textsubscript{I}]{14}
où $u^*$ est un homomorphisme d'anneaux.
Sous les mêmes hypothèses, soient $M$ un $A$-module, $N$ un $B$-module, $v:M\to N$ un homomorphisme de $A$-modules (pour la structure de $B$-module sur $N$ définie par $u:A\to B$); alors $v$ se factorise d'une seule manière
\[
  v:M\xrightarrow{i_M^S}S^{-1}M\xrightarrow{v^*}N
\]
où $v^*$ est un homomorphisme de $S^{-1}A$-modules (pour la structure de $S^{-1}A$-module sur $N$ définie par $u^*$).
\end{env}

\begin{env}[1.2.5]
\label{0.1.2.5-fr}
On définit un isomorphisme canonique $S^{-1}A\otimes_A M\xrightarrow{\sim}S^{-1}M$ de $S^{-1}A$-modules, en faisant correspondre à l'élément $(a/s)\otimes m$ l'élément $(am)/s$, l'isomorphisme réciproque appliquant $m/s$ sur $(1/s)\otimes m$.
\end{env}

\begin{env}[1.2.6]
\label{0.1.2.6-fr}
Pour tout idéal $\mathfrak{a}'$ de $S^{-1}A$, $\mathfrak{a}=(i_A^S)^{-1}(\mathfrak{a}')$ est un idéal de $A$, et $\mathfrak{a}'$ est l'idéal de $S^{-1}A$ engendré par $i_A^S(\mathfrak{a})$, qui s'identifie à $S^{-1}\mathfrak{a}$ (1.3.2).
L'application $\mathfrak{p}'\mapsto(i_A^S)^{-1}(\mathfrak{p}')$ est un isomorphisme, pour la structure d'ordre, de l'ensemble des idéaux \emph{premiers} de $S^{-1}A$ sur l'ensemble des idéaux premiers $\mathfrak{p}$ de $A$ tels que $\mathfrak{p}\cap S=\varnothing$.
En outre, les anneaux locaux $A_{\mathfrak{p}}$ et $(S^{-1}A)_{S^{-1}\mathfrak{p}}$ sont alors canoniquement (1.5.1) isomorphes.
\end{env}

\begin{env}[1.2.7]
\label{0.1.2.7-fr}
Lorsque $A$ est un anneau \emph{intègre}, dont on désigne par $K$ le corps des fractions, l'application canonique $i_A^S:A\to S^{-1}A$ est injective pour toute partie multiplicative $S$ ne contenant pas $0$, et $S^{-1}A$ s'identifie alors canoniquement à un sous-anneau de $K$ contenant $A$.
En particulier, pour tout idéal premier $\mathfrak{p}$ de $A$, $A_{\mathfrak{p}}$ est un anneau local contenant $A$, d'idéal maximal $\mathfrak{p}A_{\mathfrak{p}}$, et on a $\mathfrak{p}A_{\mathfrak{p}}\cap A=\mathfrak{p}$.
\end{env}

\begin{env}[1.2.8]
\label{0.1.2.8-fr}
Si $A$ est un anneau \emph{réduit} (1.1.1), il en est de même de $S^{-1}A$~: en effet, si $(x/s)^n=0$ pour $x\in A$, $s\in S$, cela signifie qu'il existe $s'\in S$ tel que $s'x^n=0$, d'où $(s'x)^n=0$, ce qui, par hypothèse, entraîne $s'x=0$, donc $x/s=0$.
\end{env}

\subsection{Propriétés fonctorielles}
\label{subsection:0.1.3-fr}

\begin{env}[1.3.1]
\label{0.1.3.1-fr}
Soient $M$, $N$ deux $A$-modules, $u$ un $A$-homomorphisme $M\to N$.
Si $S$ est une partie multiplicative de $A$, on définit un $S^{-1}A$-homomorphisme $S^{-1}M\to S^{-1}N$, noté $S^{-1}u$, en posant $(S^{-1}u)(m/s)=u(m)/s$; si $S^{-1}M$ et $S^{-1}N$ sont canoniquement identifiés à $S^{-1}A\otimes_A M$ et $S^{-1}A\otimes_A N$ (1.2.5), $S^{-1}u$ est identifié à $1\otimes u$.
Si $P$ est un troisième $A$-module, $v$ un $A$-homomorphisme $N\to P$, on a $S^{-1}(v\circ u)=(S^{-1}v)\circ(S^{-1}u)$; autrement dit $S^{-1}M$ est un \emph{foncteur covariant en $M$}, de la catégorie des $A$-modules dans celle des $S^{-1}A$-modules ($A$ et $S$ étant fixés).
\end{env}

\begin{env}[1.3.2]
\label{0.1.3.2-fr}
Le foncteur $S^{-1}M$ est \emph{exact}; autrement dit, si la suite
\[
  M\xrightarrow{u}N\xrightarrow{v}P
\]
est exacte, il en est de même de la suite
\[
  S^{-1}M\xrightarrow{S^{-1}u}S^{-1}N\xrightarrow{S^{-1}v}S^{-1}P.
\]
En particulier, si $u:M\to N$ est injectif (resp. surjectif), il en est de même de $S^{-1}u$;

